% !TEX root = task-6-6-2.tex
% Editable task bank. Append new BankTask blocks at the end of the document.
\documentclass[12pt,a4paper]{article}
\usepackage{physics-textbook}
\usepackage{graphicx,needspace}
\usetikzlibrary{decorations.pathmorphing}
\newif\ifShowSolutions
\ifdefined\TutorEdition
  \ShowSolutionstrue
\else
  \ShowSolutionsfalse % Change to \ShowSolutionstrue to display solutions.
\fi
\newcounter{banktask}
\newsavebox{\BankTaskBox}
\newcommand{\BankTaskContent}[3]{%
  \noindent\textbf{\thebanktask.}\quad #2\par
  \ifShowSolutions
    \Needspace{50mm}% Keep the solution heading with its opening text or diagram.
    #3\par
    \ifdefined\TutorEdition
      \ifcsname TutorSource#1\endcsname
        \nopagebreak\smallskip{\small\csname TutorSource#1\endcsname\par}%
      \fi
    \fi
  \fi
}
\newcommand{\BankTask}[3]{%
  \refstepcounter{banktask}\par\medskip
  \sbox{\BankTaskBox}{\begin{minipage}{\linewidth}%
    \BankTaskContent{#1}{#2}{#3}\end{minipage}}%
  \ifdim\dimexpr\ht\BankTaskBox+\dp\BankTaskBox\relax>\dimexpr\textheight-20mm\relax
    \Needspace{35mm}\BankTaskContent{#1}{#2}{#3}%
  \else
    \Needspace{\dimexpr\ht\BankTaskBox+\dp\BankTaskBox\relax}%
    \noindent\usebox{\BankTaskBox}\par
  \fi
  \medskip
}
\begin{document}
\BookSubsection{6.6.2}{Атом водорода}


% Source sdamgia-9513: https://phys-ege.sdamgia.ru/problem?id=9513
\BankTask{sdamgia-9513}{
Электрон в атоме водорода перешел с низкой стационарной орбиты (с меньшим номером $n$) на более высокую стационарную орбиту (с бóльшим номером $n$). Как в результате этого изменились модуль силы электрического взаимодействия электрона с ядром и полная энергия электрона?\par1. Увеличится. 2. Уменьшится. 3. Не изменится.\par\TaskChangeTable{Модуль силы электрического взаимодействия электрона с ядром}{Полная энергия электрона}
}{
\smallskip\noindent\textbf{Решение}\par
Сила электрического взаимодействия (сила Кулона) обратно пропорциональна квадрату расстояния между заряженными частицами. Поскольку электрон переходит на более высокую стационарную орбиту, увеличивается расстояние между электроном и положительно заряженным ядром, а следовательно, уменьшается модуль силы электрического взаимодействия. 
Уровни энергии электрона в атоме водорода задаются формулой $E_n=-13{,}6\,\text{эВ}/n^2$, где $n=1,2,3,\ldots$ При переходе электрона на более высокую орбиту его полная энергия увеличивается.

Ответ: $21$.
}

% Source sdamgia-10326: https://phys-ege.sdamgia.ru/problem?id=10326
\BankTask{sdamgia-10326}{
\TaskMatchingColumns{Физическая величина}{А) Энергия электрона, находящегося в атоме водорода на энергетическом уровне с номером $n$.\par Б) Энергия, которую нужно сообщить электрону в атоме водорода для того, чтобы он перешёл с $n$-го энергетического уровня на $m$-й энергетический уровень.}{Формула}{1) $E_{nm}=13{,}6\,\text{эВ}\left(\frac1{n^2}-\frac1{m^2}\right)$\par 2) $E_{nm}=13{,}6\,\text{эВ}\left(\frac1n-\frac1m\right)$\par 3) $E_n=-\frac{13{,}6\,\text{эВ}}{n^2}$\par 4) $E_n=-\frac{13{,}6\,\text{эВ}}n$}\begin{center}\begin{tabular}{|c|c|}\hline А&Б\\\hline &\\\hline\end{tabular}\end{center}
}{
\smallskip\noindent\textbf{Решение}\par
Уровни энергии электрона в атоме водорода даются формулой 

\[E_n= \frac{- 13{,}6\,\text{эВ}}{n ^2} ,\]

где $n=1,2,3,\ldots$ Чтобы атом перешел из состояния $E_{n}$ в состояние $E_{m}$ ему надо сообщить энергию 

\[E_{nm}=E_m - E_n=13{,}6\,\text{эВ} ( \frac{1}{n ^2} - \frac{1}{m ^2} ) .\]

Ответ: $31$.
}

% Source sdamgia-6204: https://phys-ege.sdamgia.ru/problem?id=6204
\BankTask{sdamgia-6204}{
В таблице приведены значения энергии для второго и четвертого энергетических уровней атома водорода.

\begin{center}\begin{tabular}{|c|c|}\hline
Номер уровня&Энергия, $10^{-19}$ Дж\\\hline
$2$&$-5{,}45$\\\hline
$4$&$-1{,}36$\\\hline
\end{tabular}\end{center}

Какой должна быть энергия фотона, при поглощении которого атом переходит со второго уровня на четвертый? (в $10^{-19}$ Дж.)
}{
\smallskip\noindent\textbf{Решение}\par
Для того, чтобы электрон перешел со второго уровня на четвертый под действием фотона, энергия фотона должна быть равна разности энергий четвертого и второго энергетических уровней: $[-1{,}36-(-5{,}45)]\cdot10^{-19}\,\text{Дж}=4{,}09\cdot10^{-19}\,\text{Дж}$.

Ответ: $4{,}09$.
}

% Source sdamgia-8951: https://phys-ege.sdamgia.ru/problem?id=8951
\BankTask{sdamgia-8951}{
Электрон в атоме водорода находится в основном (самом низком, с номером n = $1$) энергетическом состоянии. Атом поглощает фотон с импульсом $6{,}8\cdot10^{-27}$ кг·м/с. Найдите номер энергетического уровня, на который в результате этого перейдет электрон.
}{
\smallskip\noindent\textbf{Решение}\par
Уровни энергии электрона в атоме водорода задаются формулой $E_n=-\frac{13{,}6\,\text{эВ}}{n^2}$, где $n=1,2,3,\ldots$ При поглощении фотона атом переходит из основного состояния $E_1$ в некоторое состояние $E_n.$ Зная импульс фотона, можем найти его энергию:

\[E=p \cdot c ( \,\text{Дж} ) = \frac{p \cdot c}{1{,}6 \cdot 10 ^{ - 19 }} ( \,\text{эВ} ) ,\]

где $c$ — скорость света. Запишем закон сохранения энергии:

\[\frac{p \cdot c}{1{,}6 \cdot 10 ^{ - 19 }} =E_n - E_1=13{,}6 \cdot ( 1 - \frac{1}{n ^2} ) \Rightarrow n ^2 = \frac{1}{1 - \frac{p \cdot c}{1{,}6 \cdot 10 ^{ - 19 } \cdot 13{,}6}} =16.\]

Отсюда находим, что электрон перейдет на уровень с номером $n=4$.

Ответ: $4$.
}

% Source sdamgia-10650: https://phys-ege.sdamgia.ru/problem?id=10650
\BankTask{sdamgia-10650}{
Электрон в атоме водорода переходит с энергетического уровня с номером n = $2$ на энергетический уровень с n = $1$. Чему равен модуль импульса испущенного при этом фотона? (В $10^{-29}\,\text{кг}\cdot\text{м/с}$; округлите до целого числа.)
}{
\smallskip\noindent\textbf{Решение}\par
Уровни энергии электрона в атоме водорода задаются формулой $E_n=-13{,}6\,\text{эВ}/n^2$, где $n=1,2,3,\ldots$ При переходе атома из состояния $E_{n}$ в состояние $E_{m}$ атом испускает фотон с энергией $E=E_n - E_m.$ При переходе со второго на первый энергетический уровень атом испустит фотон с энергией

\[E=E_2 - E_1=13{,}6\,\text{эВ}\cdot(1-\frac14) =10{,}2\,\text{эВ}.\]

Импульс фотона связан с его энергией как

\[p=\frac{E}{c}=\frac{10{,}2\cdot1{,}6\cdot10^{-19}\,\text{Дж}}{3\cdot10^8\,\text{м/с}}=544\cdot10^{-29}\,\text{кг}\cdot\text{м/с}.\]

Ответ: $544$.
}

% Source sdamgia-4758: https://phys-ege.sdamgia.ru/problem?id=4758
\BankTask{sdamgia-4758}{
Уровни энергии электрона в атоме водорода задаются формулой $E_n= - \frac{13{,}6}{n ^2} \,\text{эВ},$ где $n = 1, 2, 3, ...$ . При переходе атома из состояния $E_2$ в состояние $E_1$ атом испускает фотон. Попав на поверхность фотокатода, фотон выбивает фотоэлектрон. Длина волны света, соответствующая красной границе фотоэффекта для материала поверхности фотокатода, 	$\lambda_{\text{кр}} = 300\,\text{нм}.$ Чему равна максимально возможная кинетическая энергия фотоэлектрона?
}{
\smallskip\noindent\textbf{Решение}\par
Согласно постулатам Бора, свет излучается при переходе атома на более низкие уровни энергии, при этом фотоны несут энергию, равную разности энергий начального и конечного состояний. Таким образом, испущенный фотон имел энергию

\[h\nu=E_2 - E_1= - 13{,}6 \,\text{эВ} ( \frac{1}{2 ^2} - \frac{1}{1 ^2} ) = 10{,}2\,\text{эВ}.\]

Согласно уравнению фотоэффекта, максимальная кинетическая энергия вылетающих фотоэлектронов связана с энергией фотона и работой выхода соотношением

\[h\nu=A + E_{\,\text{к}}.\]

Работа выхода связана с длиной волны красной границы соотношением: 

\[A= \frac{hc}{\lambda_{\,\text{кр}}}\]

Таким образом, максимально возможная кинетическая энергия фотоэлектрон равна

\[E_{\,\text{к}}=h\nu - \frac{hc}{\lambda_{\,\text{кр}}} = 10{,}2 \,\text{эВ} - \frac{4{,}1 \cdot 10 ^{ - 15 } \,\text{эВ} \cdot \,\text{с} \cdot 3 \cdot 10 ^{8} \,\text{м/с}}{3 \cdot 10 ^{ - 7 } \,\text{м}} = 6{,}1 \,\text{эВ}.\]

Ответ: $E_{\,\text{к}} = 6{,}1\,\text{эВ}$
}

% Source sdamgia-5388: https://phys-ege.sdamgia.ru/problem?id=5388
\BankTask{sdamgia-5388}{
Электрон, имеющий импульс $p=2\cdot10^{-24}\,\text{кг}\cdot\text{м/с}$ сталкивается с покоящимся протоном, образуя атом водорода в состоянии с энергией $E_{n}$ ($n = 2$). В процессе образования атома излучается фотон. Найдите частоту $\nu$ этого фотона, пренебрегая кинетической энергией атома. Уровни энергии электрона в атоме водорода задаются формулой $E_n= - \frac{13{,}6}{n ^2} \,\text{эВ},$ где $n=1,2,3,... .$
}{
\smallskip\noindent\textbf{Решение}\par
Поскольку кинетической энергией атома пренебрегаем, то вся кинетическая энергия электрона $E_0= \frac{p ^2}{2m_e}$ перейдет в энергию атома $E_2= - \frac{13{,}6}{2 ^2} \,\text{эВ}$ и энергию фотона.

Таким образом, по закону сохранения энергии энергия фотона равна 

\[E_{\,\text{фот}} =E_0 - E_2= \frac{p ^2}{2m_e} + \frac{13{,}6\,\text{эВ}}{2 ^2} =h\nu.\]

Отсюда находим частоту фотона:

\[\nu= \frac{p ^2 / ( 2m_e ) + (13{,}6\,\text{эВ})/2 ^2}{h}\]

 Тогда 

$\nu= \frac{4 \cdot 10 ^{ - 48 } \,\text{кг} \cdot \,\text{Дж}/ ( 2 \cdot 9{,}1 \cdot 10 ^{ - 31 } \,\text{кг} ) + 13{,}6/4 \cdot 1{,}6 \cdot 10 ^{ - 19 } \,\text{Дж/эВ} \cdot \,\text{эВ}}{6{,}62 \cdot 10 ^{ - 34 } \,\text{Дж} \cdot \,\text{с}} =$

$=4{,}14 \cdot 10 ^{ 15 } \,\text{Гц}\approx 4{,}1 \cdot 10 ^{ 15 } \,\text{Гц}.$

Ответ: $\nu=4{,}1 \cdot 10 ^{ 15 } \,\text{Гц}.$
}

% Source sdamgia-6513: https://phys-ege.sdamgia.ru/problem?id=6513
\BankTask{sdamgia-6513}{
Значения энергии электрона в атоме водорода задаются формулой $E_n= \frac{- 13{,}6\,\text{эВ}}{n ^2} ,n=1,2,3,....$ При переходе с верхнего уровня энергии на нижний атом излучает фотон. Переходы с верхних уровней на уровень с $n=1$ образуют серию Лаймана; на уровень с $n=2$  — серию Бальмера; на уровень с $n = 3$  — серию Пашена и так далее. Найдите отношение $\beta$ минимальной частоты фотона в серии Бальмера к максимальной частоте фотона в серии Пашена.
}{
\smallskip\noindent\textbf{Решение}\par
Частота фотона связана с его энергией равенством $h\nu=E,$ где $h$  — постоянная Планка. В серии Бальмера энергия фотона равна $E_n - E_2,$ где $n=3, 4, ....$ Аналогично в серии Пашена энергия фотона равна $E_n - E_3,$ где $n=4, 5, ....$ Минимальной частота фотона в серии Бальмера будет при условии перехода с $3$-го уровня, максимальной частота фотона в серии Пашена будет при переходе с самого высокого ($n = \infty$) уровня. Поэтому 

\[\beta = \frac{E_3 - E_2}{E_{\infty} - E_3} = \frac{\frac{1}{2 ^2} - \frac{1}{3 ^2}}{\frac{1}{3 ^2} - 0} =1{,}25.\]

Ответ: $1{,}25$.
}

\end{document}
